Large-eddy simulation (LES) is the de facto standard for wind-farm aerodynamics
because it directly resolves energy-carrying turbulence in the atmospheric
boundary layer rather than modeling it. LESGO is a widely used open-source
solver for this regime built on pseudo-spectral methods, which carry minimal
numerical dissipation and deliver high accuracy per grid point. LESGO's CPU
implementation, however, suffers from poor parallel scalability, severely
constraining the grid resolution and physical integration times of
production-scale wind-farm simulations. To enable significantly finer resolutions and longer time scales, we
present the porting of LESGO to GPGPUs. For this class of pseudo-spectral
multi-physics solvers with recurring host--device coupling, data movement rather
than kernel throughput governs overall performance. We analyze all host--device
coupling points across the solver and classify them into four treatments. On
this foundation, we propose several optimizations, including chunk pipelining
across serial rank chains, batching many small work items, and heterogeneous
CPU--GPU execution.  Evaluation on a 60-turbine wind farm discretized into 604 million cells
demonstrates a 40$\times$ speedup over the best-performing CPU baseline at
equal node count, rising to 90$\times$ on 64 GPUs, and strong scaling to
128 GPUs on a 1.36 billion cell grid.
